% ═══ Results: Boolean Logic over Conformational State ═══
\section{Boolean Logic over Conformational State}
\label{sec:res-contact-logic}

The results so far apply Boolean minimisation to \emph{sequence}: an alignment
column pair becomes a Karnaugh map, and its prime implicants describe which
residue combinations coevolution permits.  Section~\ref{sec:res-ceiling} shows
that this substrate is close to exhausted for contact prediction.  Here we
change the substrate rather than the method.

The motivation is that Boolean minimisation is productive elsewhere in biology
precisely where the object is a \emph{logical relation among variables observed
across many states}---gene-regulatory inference being the established case, in
which binarised expression time series yield logic rules via
Quine--McCluskey~\cite{kocevar2024sailor}.  A static residue string is not such
an object.  A structure-based model (SBM) trajectory~\cite{clementi2000sbm}
is: every native contact is either formed or broken in every frame, so the
trajectory is a binary state matrix and one may ask whether a contact's state
is a Boolean function of a few other contacts' states.

\subsection{Controlling the folding coordinate}
\label{sec:cl-control}

Every pair of native contacts is correlated for a trivial reason: when the
protein is folded most contacts are formed, when unfolded most are broken.  Any
naive coupling statistic measures this global mode and nothing else.  Three
devices control for it, all applied throughout.

First, $Q$, the fraction of native contacts formed, is quantile-binned into two
bits and supplied to the circuit as an \emph{explicit input}, so that a rule
reads ``at this folding level, contact $c$ is formed iff \dots'' rather than
``contact $c$ is formed when the protein is folded''.  Second, every rule is
scored against a \emph{$Q$-only baseline} built by the identical procedure from
the $Q$ bits alone; the reported quantity is always the gain over that
baseline, so a rule that merely restates the global mode scores zero.  Third, a
\emph{$Q$-stratified null} resamples contact columns within $Q$ strata,
preserving each contact's dependence on the folding coordinate and its marginal
frequency while destroying any specific contact-to-contact coupling.

Predictors are ranked by \emph{conditional} mutual information
$I(c; x \mid Q\text{-bin})$; plain mutual information would rank candidates by
how strongly each tracks the confounder.  Predictor selection is repeated
inside every cross-validation fold, so the \emph{choice} of predictors is held
out and not only the fitted rule.  Folds are contiguous frame blocks rather
than random splits.

\subsection{Coupling beyond the folding coordinate is present}

Across 480 proteins and 18{,}714 extracted rules, every circuit
runtime-verified sound and complete, the mean gain over the $Q$-only baseline
is $-0.0054$ while the $Q$-stratified null gives $-0.0248$.  The real
configuration exceeds its own null in \textbf{463 of 480 proteins}
(Wilcoxon $p = 4.3\times10^{-78}$).

The sign deserves comment.  On held-out frames the contact-based circuit is
marginally \emph{worse} than the two-bit $Q$ model, yet decisively better than
its own null.  Varying the predictor count identifies the cause
(Table~\ref{tab:cl-predcount}): the absolute gain degrades monotonically as the
truth table grows, while the real-minus-null difference stays flat at
$\approx +0.024$ and equally significant at every size.  That is a
variance-limited estimator, not absent signal---with 100 frames there are
$\approx 2.5$ frames per cell at three predictors.

\begin{table}[htbp]
\centering
\caption{Predictor count against estimator variance (60 proteins, 25 targets
each).  Absolute gain falls with table size; the real-minus-null difference
does not.}
\label{tab:cl-predcount}
\begin{tabular}{rrrrrl}
\toprule
Predictors & Cells & Gain vs.\ $Q$-only & Null gain & Real $-$ null & Real $>$ null \\
\midrule
1 & 8  & $+0.0033$ & $-0.0195$ & $+0.0228$ & 52/55 ($p = 2.8\times10^{-10}$) \\
2 & 16 & $-0.0002$ & $-0.0259$ & $+0.0257$ & 50/55 ($p = 6.0\times10^{-10}$) \\
3 & 32 & $-0.0023$ & $-0.0260$ & $+0.0237$ & 51/55 ($p = 3.5\times10^{-10}$) \\
\bottomrule
\end{tabular}
\end{table}

\subsection{The learned rules are spatially local}
\label{sec:cl-local}

The substantive finding is what the rules connect.  For each rule we compare
the contact-map $L_1$ distance $|i-k| + |j-l|$ between the target contact and
its selected predictors against predictors drawn at random from the
\emph{same protein's} varying contacts, so that protein size and contact
density cancel.

\begin{table}[htbp]
\centering
\caption{Selected predictors versus matched-random predictors, 480 proteins.
Adjacency is $L_1 \leq 2$ on the contact map.}
\label{tab:cl-locality}
\begin{tabular}{lrr}
\toprule
 & Observed & Matched random \\
\midrule
Fraction of predictors adjacent      & \textbf{0.2127} & 0.0170 \\
Enrichment                           & \multicolumn{2}{c}{\textbf{12.5$\times$}} \\
Proteins with observed $>$ random    & \multicolumn{2}{c}{\textbf{478 / 480}} \\
Wilcoxon                             & \multicolumn{2}{c}{$p = 3.3\times10^{-80}$} \\
Median contact-map distance          & \textbf{61.5} & 107.5 \\
Predictors sharing a residue         & \multicolumn{2}{c}{\textbf{27.2\%}} \\
\bottomrule
\end{tabular}
\end{table}

Conditional mutual information, with the global folding mode controlled out,
selects the \emph{neighbouring cells of the contact map}: adjacent rungs of the
same $\beta$-ladder, successive turns packed against the same partner.  More
than a quarter of selected predictors share a residue with their target.  This
is local cooperativity---the foldon picture~\cite{maity2005foldons}---recovered
from a discrete-logic direction without secondary structure being supplied.
Individual rules can be strong even where the population mean gain is not:

\begin{quote}\small
\texttt{1zjy}\quad contact $(1,28)$: accuracy $0.900$ vs.\ $Q$-only $0.480$, MCC $0.80$;\\
\phantom{\texttt{1zjy}}\quad predictors $(1,29), (2,28), (2,29)$, contact-map $L_1 = 1, 1, 2$.

\texttt{1kms}\quad contact $(2,129)$: accuracy $0.790$ vs.\ $0.380$, MCC $0.62$;\\
\phantom{\texttt{1kms}}\quad rule $c(2,129) = c(2,127) \wedge c(1,129)$, $L_1 = 2, 1$.
\end{quote}

\subsection{Longer trajectories confirm the variance diagnosis}
\label{sec:cl-dense}

The diagnosis in Table~\ref{tab:cl-predcount} makes a testable prediction:
more statistically independent configurations should flip the absolute gain
positive.  We re-simulated ten proteins for four times longer with twenty times
denser output ($2\times10^{6}$ steps saved every 250, against $5\times10^{5}$
every 5000) and repeated the identical analysis in three arms
(Table~\ref{tab:cl-dense}).

\begin{table}[htbp]
\centering
\caption{Three arms on the same ten proteins.  The middle arm decimates the
long run back to the \emph{original} 5000-step spacing, so its time resolution
is unchanged and only the number of independent configurations differs.}
\label{tab:cl-dense}
\begin{tabular}{lrrrrr}
\toprule
Condition & Frames & Spacing & Gain vs.\ $Q$-only & Null & Fraction $>0$ \\
\midrule
Original     & 100   & 5000 & $+0.0164$ & $-0.0549$ & 0.90 \\
Long-sparse  & 400   & 5000 & $\mathbf{+0.0628}$ & $-0.0299$ & \textbf{1.00} \\
Long-dense   & 8000  & 250  & $\mathbf{+0.0842}$ & $-0.0022$ & \textbf{1.00} \\
\bottomrule
\end{tabular}
\end{table}

Long-sparse is the informative arm: identical time resolution to the original,
four times the independent configurations, and the gain rises $3.8\times$ and
becomes positive in every protein ($+0.0464$, 10/10, $p = 0.00195$).  The
constraint was sampling, not method.  These ten proteins were selected as small
and strongly fluctuating, so their absolute levels are optimistic relative to
the 480-protein population; the paired comparisons are unaffected.

\subsection{Dense sampling licenses a kinetic reading}
\label{sec:cl-kinetic}

In the original trajectories the lag-1 autocorrelation of contact state
averages $0.009$ ($Q$: $0.130$): the 5000-step save interval exceeds the
correlation time, the frames are near-independent equilibrium samples, and no
$t \rightarrow t+1$ statement is licensed.  We measured this per protein rather
than assuming it, and made no kinetic claim on that data.

At 250-step spacing the contact lag-1 autocorrelation rises to $0.539$.  A
lagged analysis---predicting the target $k$ frames \emph{ahead} of the inputs,
with predictor selection lagged identically---therefore becomes meaningful
(Table~\ref{tab:cl-lag}).

\begin{table}[htbp]
\centering
\caption{Predicting a contact ahead of its inputs, ten densely sampled
proteins.}
\label{tab:cl-lag}
\begin{tabular}{rlrrl}
\toprule
Lag & Horizon & Gain vs.\ $Q$-only & Null & Real $>$ null \\
\midrule
0 & same frame  & $+0.0842$ & $-0.0022$ & 10/10 ($p = 1.95\times10^{-3}$) \\
1 & 250 steps   & $+0.0593$ & $-0.0023$ & 10/10 ($p = 1.95\times10^{-3}$) \\
2 & 500 steps   & $+0.0431$ & $-0.0024$ & 10/10 ($p = 1.95\times10^{-3}$) \\
4 & 1000 steps  & $+0.0288$ & $-0.0021$ & 10/10 ($p = 1.95\times10^{-3}$) \\
\bottomrule
\end{tabular}
\end{table}

Predictive power decays roughly exponentially, with a relaxation of
$\approx 900$ steps, but remains positive in every protein a thousand
integration steps ahead.  The state of three neighbouring contacts now carries
information about a target contact's future state beyond what the folding level
supplies---a statement about mechanism rather than only about the equilibrium
ensemble, and the first such statement this framework has been entitled to
make.
